% !TeX root = ../main.tex

\chapter{Implementation}\label{chapter:implementation}

\section{Repository layout}

The InsertGS code is organized as a single Python package with three
external entry points:

\begin{description}
  \item[\texttt{main.py}] Unified CLI and Gradio dashboard. With no
    arguments it serves the dashboard on port 7860; with a subcommand
    (\texttt{list}, \texttt{info}, \texttt{batch}, \texttt{bg},
    \texttt{compose}) it runs the corresponding pipeline step.
  \item[\texttt{tools/serve\_spz.py}] A lightweight web server that hosts
    \ac{SPZ} scenes together with the SuperSplat editor and viewer, so
    captured rooms can be inspected directly in the browser.
  \item[\texttt{tools/spz\_to\_ply.py}] An \ac{SPZ} $\to$ 3DGS \ac{PLY}
    decoder, used both as a standalone CLI and as an imported helper.
\end{description}

\section{Pipeline modules}

The pipeline stages from \Cref{chapter:method} live under
\texttt{src/scripts/}: cubemap extraction
(\texttt{gs\_to\_cubemap\_gsplat.py}), Blender PBR rendering
(\texttt{render\_chair.py} and a Cycles wrapper), and the compositing pass.
Inputs --- room PLY files and object meshes --- are gathered under
\texttt{src/inputs/rooms/} and \texttt{src/inputs/objects/}, while each
experiment writes its outputs into \texttt{src/cases/}\textit{name}\,/.

\section{Environment}

InsertGS runs in a dedicated conda environment named \texttt{InsertGS}.
The Python entry point auto-injects \texttt{CUDA\_HOME} pointing at the
conda prefix so that \texttt{gsplat} compiles its CUDA extensions on
demand. Blender is invoked out-of-process from
\texttt{/home/ai/bin/blender}, which decouples the rendering Python from
the rest of the pipeline.

\section{Web viewer}

A thin HTML/JS page (\texttt{viewer.html}) acts as a scene picker for the
\texttt{serve\_spz.py} server. Selecting a scene opens it in the bundled
SuperSplat viewer; SPZ scenes are decoded on the fly when needed. This is
used both during dataset curation and to share intermediate results
without redistributing the heavyweight PLY files.
